% ============================================================================
\section{Búsqueda Tabú}\label{sec:tabu}
% ============================================================================

\subsection{Idea}

La Búsqueda Tabú (\emph{Tabu Search}, Glover, 1986; Hansen, 1986) es una búsqueda local que
\textbf{no se detiene} al llegar a un óptimo local. En cada iteración ejecuta el mejor
movimiento disponible \emph{aunque empeore} la solución. El problema evidente es que, tras un
movimiento que empeora, el mejor movimiento siguiente sería deshacerlo, y el algoritmo
entraría en un ciclo sin fin entre dos soluciones. La solución es dotarlo de
\textbf{memoria}:

\begin{itemize}
  \item \textbf{Memoria a corto plazo (lista tabú).} Los movimientos hechos recientemente
        quedan prohibidos (son ``tabú'') durante un cierto número de iteraciones, que se
        llama \emph{duración} tabú. Así no se puede volver atrás enseguida y la búsqueda se
        ve obligada a alejarse del óptimo local.
  \item \textbf{Función de aspiración.} Una prohibición puede impedir un movimiento
        excelente. Por eso se permite saltarse el tabú si el movimiento cumple un criterio
        de aspiración; lo habitual (dicen los apuntes) es aceptarlo \emph{si mejora la mejor
        solución conocida hasta el momento}.
  \item \textbf{Memoria a largo plazo.} Se registra la \emph{frecuencia} con que se ha hecho
        cada movimiento desde el principio, y se usa para \emph{diversificar} (penalizar lo
        que se ha hecho muchas veces, para visitar zonas nuevas) o para \emph{intensificar}
        (volver a lo que ha dado buenos resultados). Según los apuntes, esto se hace
        habitualmente añadiendo a la función objetivo una penalización proporcional a la
        frecuencia del movimiento.
\end{itemize}

A diferencia del temple, la búsqueda tabú es esencialmente determinista: no decide al azar,
sino que usa la historia de la búsqueda. Un tutorial clásico y de acceso libre
es~\cite{glover}.

\subsection{Esquema de los apuntes}

\begin{apuntes}[Esquema de los apuntes: esquema de un algoritmo de búsqueda tabú]
\begin{tabbing}
\quad\=\quad\=\kill
Inicialización: \ $s\in X$; \ $s^\ast=s$; \ $k=0$; \ $T=\emptyset$; \ $A(z)=z$\\
Proceso:\\
\> \textbf{do}\\
\>\> $M(s,k)=\{\,m\in M \;/\; m\notin T \ \vee\ f(s\oplus m)<A(f(s))\,\}$ \quad (movimientos válidos)\\
\>\> $V^\ast=\{\,s'\in X \;/\; s'=s\oplus m,\ m\in M(s,k)\,\}$ \quad (conjunto de vecinos)\\
\>\> Se elige la mejor solución $s'\in V^\ast$\\
\>\> If $(f(s')<f(s^\ast))$ \ $s^\ast=s'$\\
\>\> Actualizar $T$ y $A(\cdot)$\\
\>\> $k=k+1$; \ $s=s'$\\
\> \textbf{while} (condición de parada)
\end{tabbing}
\end{apuntes}

En palabras: $s\oplus m$ es la solución que resulta de aplicar el movimiento $m$ a $s$. En
cada iteración se forma el conjunto $M(s,k)$ de movimientos \emph{válidos} (los que no están
en la lista tabú $T$, más los que sí están pero cumplen la aspiración), se ejecuta el mejor
de ellos y se actualizan la mejor solución $s^\ast$ y la lista tabú.

\subsection{Adaptación al problema de cubrimiento}

\paragraph{Movimientos.} Se necesitan movimientos que permitan tanto mejorar como empeorar de
forma controlada, manteniendo la factibilidad.

\begin{decision}[Decisión: dos tipos de movimiento simétricos]
\begin{itemize}[itemsep=1pt]
  \item \textbf{Entrada} de una columna $k$: entra $k$ y salen las columnas que $k$ vuelve
        redundantes. Es el movimiento de $\D3$. Valor: $v=c_k-\sum c_{\text{salen}}$.
  \item \textbf{Salida} de una columna $j$: sale $j$ y las filas que solo cubría ella se
        vuelven a cubrir con el voraz del cociente (entran una o varias columnas). Valor:
        $v=\sum c_{\text{entran}}-c_j$.
\end{itemize}
La entrada cambia ``varias por una'' y la salida ``una por varias''. Hacen falta las dos: con
solo entradas la búsqueda únicamente puede reducir el número de columnas y se queda sin
movimientos útiles en cuanto llega a un óptimo local de $\D3$; las salidas le permiten
desandar el camino por otro sitio. En cada iteración se valoran \emph{todos} los movimientos
de ambos tipos y se ejecuta el de menor valor admisible, sea negativo (mejora) o positivo
(empeora).
\end{decision}

\paragraph{Qué se declara tabú.} No se guardan soluciones completas (compararlas sería
carísimo), sino \emph{atributos} de los movimientos: las columnas que han entrado o salido.

\begin{decision}[Decisión: dos listas tabú, una de entrada y otra de salida]
\begin{itemize}[itemsep=1pt]
  \item Una columna que \textbf{sale} no puede volver a \textbf{entrar} durante un tiempo.
        Impide deshacer el movimiento.
  \item Una columna que \textbf{entra} no puede \textbf{salir} durante un tiempo. Impide que
        la expulsen nada más llegar, antes de que la búsqueda haya podido aprovecharla.
\end{itemize}
Cada lista tiene su propia duración (\cod{tent} y \cod{tsal}) porque no son simétricas: hay
casi 9\,900 columnas fuera de la solución y solo unas 130 dentro. Prohibir la entrada de 20
columnas apenas restringe; prohibir la salida de 20 columnas bloquea una parte apreciable de
la solución. Por eso cabe esperar que convenga $\cod{tsal}<\cod{tent}$; la calibración lo
comprueba (sección~\ref{sec:cal-tabu}).
\end{decision}

\begin{decision}[Decisión: cómo se implementa la lista tabú]
No hay una ``lista'' propiamente dicha. Para cada columna $j$ se guarda el número de la
iteración hasta la cual está prohibida: \cod{tabu\_ent[j]} y \cod{tabu\_sal[j]}. Comprobar si
un movimiento es tabú es una comparación (\cod{tabu\_ent[k] > iter}) y las prohibiciones
caducan solas, sin tener que borrar nada. Es equivalente a la ``estructura de datos tabú'' de
los apuntes, que guarda cuántas iteraciones de prohibición le quedan a cada movimiento, pero
evita tener que restar uno a todas las casillas en cada iteración.
\end{decision}

\begin{decision}[Decisión: duración aleatoria]
La duración de cada prohibición se sortea entre \cod{ten} y $2\cdot\cod{ten}$. Con una
duración fija, la búsqueda puede caer en ciclos de longitud mayor que la duración; variarla
al azar los rompe. Es la recomendación de~\cite[sec.~3.2.3]{glover}.
\end{decision}

\paragraph{Función de aspiración.} Los apuntes dan dos ejemplos, $A(z)=z$ y
$A(z)=C(x^\ast)$, y añaden que ``es habitual solo aceptar un movimiento tabú si mejora la
mejor solución conocida hasta el momento''. Se usa esta segunda, $A=f(s^\ast)$: un
movimiento tabú se admite si lleva a una solución mejor que la mejor conocida. No puede
producir ciclos (una solución mejor que todas las anteriores es, por fuerza, nueva), cosa que
$A(z)=z$ no garantiza.

\paragraph{Memoria a largo plazo: diversificación.} \cod{frec[j]} cuenta cuántas veces ha
entrado la columna $j$ en la solución. Al comparar movimientos no se usa el valor $v$ sino el
valor penalizado
\[
   v' \;=\; v \;+\; \beta\cdot\sum_{j\ \text{que entra}} \cod{frec}[j] ,
\]
de modo que, a igualdad de valor, se prefieren las columnas que se han usado poco. Es la
misma idea del ejemplo de los apuntes~\cite{yanez}, donde al valor de cada movimiento se le
resta su frecuencia absoluta (allí el problema es de máximo; aquí, de mínimo, se suma). El
peso $\beta$ es necesario porque nuestros valores son costes reales y las frecuencias son
números enteros: sin él no serían comparables. Con $\beta=0$ no hay memoria a largo plazo.

\paragraph{Intensificación.} Si pasan \cod{reinicio} iteraciones sin mejorar $s^\ast$, la
búsqueda vuelve a $s^\ast$ con las listas tabú vacías. Las frecuencias \emph{no} se borran,
así que la búsqueda no repetirá el mismo camino que la vez anterior.

\paragraph{Limpieza.} Tras un movimiento pueden quedar columnas redundantes. Se eliminan las
que no estén protegidas por el tabú de salida, y cada una queda con prohibición de volver a
entrar.

\subsection{Parámetros}

\begin{center}
\begin{tabular}{@{}lp{11.3cm}@{}}
\toprule
Nombre & Significado y efecto \\ \midrule
\cod{tent} & Duración tabú de \emph{entrada} (para las columnas que salen). Corta: riesgo de
             ciclos. Larga: se prohíben demasiadas columnas buenas. \\
\cod{tsal} & Duración tabú de \emph{salida} (para las columnas que entran). \\
\cod{beta} & Peso $\beta$ de la penalización por frecuencia. Mayor: más diversificación. \\
\cod{reinicio} & Iteraciones sin mejora antes de volver a $s^\ast$. Menor: más
             intensificación. Un valor enorme equivale a no reiniciar nunca. \\
\bottomrule
\end{tabular}
\end{center}

\subsection{Recorrido por el código}

Cada iteración empieza valorando todas las entradas. Obsérvese la condición de admisibilidad
(tabú y aspiración) y la penalización por frecuencia, que corresponden a la definición de
$M(s,k)$ del esquema:

\fragmento{tabu.c}{195}{208}

Después se valoran las salidas con \cod{valor\_salida}, que simula la reparación voraz
mirando solo las filas únicas de la columna que sale (son muy pocas), sin tocar la solución.
Elegido el mejor movimiento, se ejecuta y se actualizan las dos listas tabú y las
frecuencias (``Actualizar $T$'' en el esquema):

\fragmento{tabu.c}{245}{262}

Por último se actualiza la mejor solución y, si procede, se intensifica:

\fragmento{tabu.c}{269}{287}

\begin{idea}[Coste de una iteración]
Una iteración tabú valora casi 9\,900 entradas y unas 130 salidas, lo que supone recorrer
los unos de la matriz. El programa hace unas 950 iteraciones por segundo. Es decir, en el
mismo tiempo la búsqueda tabú da unas 90 veces menos pasos que el temple, pero cada paso es
el \emph{mejor} de unos 10\,000 movimientos y no uno al azar. Son dos formas opuestas de
gastar el tiempo de cálculo.
\end{idea}
